% ===========================================================================
% P1.10 NHÁP - các phần mới/viết lại cho main_vi.tex (khung A). Bản dịch của
% p1_10_sections_en.tex, CÙNG số liệu (nguồn: outputs/phase5/plan_v3/).
% Viết 2026-10-07, CHƯA gộp vào main_vi.tex: chờ quyết định C2 (tên bài),
% C4 (tạp chí), C5 (force_periodic). Dùng macro của main_vi.tex.
% ===========================================================================

% --- THAY tóm tắt -----------------------------------------------------------
\begin{abstract}
Thiết kế ô cơ sở vật liệu meta dựa trên gradient, dù bằng tối ưu tô-pô hay
bằng tìm kiếm trong không gian ẩn của mô hình sinh, được đánh giá bằng phân
tích phần tử hữu hạn (FE) trên thiết kế đã nhị phân hóa, nhưng thường lại tối
ưu một trường mật độ liên tục. Chúng tôi chỉ ra rằng sự lệch này bị khai thác
một cách có hệ thống. Tinh chỉnh mã ẩn của một bộ mã hóa tự động biến phân có
điều kiện (cVAE) theo độ nhạy đồng nhất hóa chính xác đưa hàm mục tiêu về
$\sim\!10^{-10}$ trong khi sai số kiểm chứng vẫn lớn hơn sáu bậc, và tối ưu
tô-pô dựa trên mật độ với Heaviside continuation tới $\beta=64$ cũng có cùng
khoảng lệch (trung vị $2\times10^{-7}$ trước so với $3{,}6\times10^{-2}$ sau
khi lấy ngưỡng). Đưa các toán tử kiểm chứng vào hàm mục tiêu tinh chỉnh nâng
$\RR$ kiểm chứng FE của hệ số Poisson $\nuot$ cho ô auxetic hai chiều từ 0,874
lên 0,998 với mẫu đơn lẻ (0,831 lên 0,990 trên tập mục tiêu thứ hai, không
giao nhau) và từ 0,983--0,990 lên 0,995--0,998 sau chọn best-of-30 trên ba
tập. So với tối ưu tô-pô chạy từ đầu với nhiều điểm khởi tạo, quy trình sinh
đạt cùng độ chính xác với số lần giải FE ít hơn khoảng tám lần, và trên một
tập kiểm định đăng ký trước, sai số kết hợp thấp hơn 27\% đối với mục tiêu
dị hướng vừa phải. Phương pháp thất bại với mục tiêu dị hướng mạnh vốn thưa
trong dữ liệu huấn luyện, và tối ưu tô-pô vẫn vượt trội về khả năng chế tạo
cũng như với mục tiêu ngoài dải huấn luyện.
\end{abstract}

% --- Results MỚI ------------------------------------------------------------
\subsection{So sánh với tối ưu tô-pô chạy từ đầu}

Đối thủ đúng nghĩa của một quy trình sinh không phải là tra bảng mà là tối
ưu tô-pô chạy trực tiếp cho mục tiêu. Với mỗi mục tiêu, chúng tôi giải bài
toán đồng nhất hóa ngược kinh điển
$\min\tfrac12\lVert\boldsymbol\nu-\boldsymbol\nu^{*}\rVert^2$ bằng phương pháp
tiệm cận di động (MMA), có cận thể tích và sàn độ cứng, Heaviside continuation
tới $\beta=64$ và bốn tô-pô khởi tạo lấy từ dữ liệu huấn luyện (Phương pháp).
Mỗi lần chạy được kiểm chứng y hệt một thiết kế sinh ra, và lời giải tốt nhất
trong bốn được chọn theo sai số kiểm chứng. Baseline này dùng trung bình 893
lần giải FE cho mỗi mục tiêu, so với 107 của best-of-30 kèm tinh chỉnh nhận
biết nhị phân hóa có bảo vệ.

Tối ưu tô-pô có cùng khoảng lệch mục tiêu--kiểm chứng như tinh chỉnh không
gian ẩn. Cuối continuation, sai số trung vị trên trường mật độ đã chiếu là
$2\times10^{-7}$, còn sai số kiểm chứng trung vị sau khi lấy ngưỡng là
$3{,}6\times10^{-2}$ mỗi lần chạy, và ở 57\% trong 300 mục tiêu, tô-pô khởi
tạo có sai số trên trường chiếu thấp nhất không phải là tô-pô có sai số kiểm
chứng thấp nhất. Do đó việc chọn theo sai số kiểm chứng là cần thiết cho cả
hai phương pháp, củng cố luận điểm trung tâm của bài vượt ra ngoài mô hình
sinh.

Bảng~\ref{tab:simp} so sánh hai phương pháp trên ba tập 100 mục tiêu giữ lại,
không giao nhau. Trên cả ba tập, quy trình sinh ngang bằng tối ưu tô-pô hội
tụ về $\nuot$ ($\RR=0{,}994$--$0{,}998$ so với $0{,}996$--$0{,}998$) với số lần
giải FE ít hơn khoảng tám lần. Trên toàn bộ mục tiêu, sai số kết hợp
$|\Delta\nuot|+|\Delta\nuto|$ không thấp hơn một cách nhất quán, do một kiểu
thất bại cụ thể. Mọi ô bị đứt mà cVAE sinh ra (một mục tiêu ở tập thứ nhất,
hai ở tập thứ hai) đều thuộc mục tiêu có tỉ số dị hướng
$r=\max(\nuto^{*}/\nuot^{*},\,\nuot^{*}/\nuto^{*})\ge10$, tương ứng tỉ số độ
cứng $E_2/E_1\ge10$ với ô trực hướng. Với các mục tiêu này cả 30 ứng viên
đều bị đứt, nên không thể sửa bằng cách chọn; chỉ 2,8\% mẫu huấn luyện có
$r>10$. Sau khi quan sát điều này ở hai tập đầu, chúng tôi đăng ký trước một
kiểm định phân tầng trên tập thứ ba trước khi chạy: với mục tiêu $r<10$, quy
trình sinh phải có sai số kết hợp và sai số $\nuot$ thấp hơn tối ưu tô-pô, với
CI 95\% ghép cặp không chứa 0. Tiêu chí sai số kết hợp đạt (thấp hơn 27\%, CI
8--42\%); tiêu chí $\nuot$ không đạt (thấp hơn 19\%, CI $-10$ đến 40\%). Kiểm định
này dùng quy trình có ép biên tuần hoàn. Sau khi bỏ bước này vì lý do khái
niệm (Phương pháp), cùng phép so sánh nghiêng về quy trình sinh ở cả hai tiêu
chí trên cả ba tập (Bảng~\ref{tab:simp}); vì phiên bản này không được đăng ký
trước, chúng tôi báo nó như phân tích độ nhạy và vẫn chỉ khẳng định ngang
bằng về $\nuot$ và lợi thế về sai số kết hợp ở dị hướng vừa phải, không khẳng
định vượt trội về $\nuot$. Chênh lệch
$\nuot$ giữa hai phương pháp ($\sim\!10^{-3}$) dù sao cũng nhỏ hơn một bậc so
với sai số rời rạc hóa của mô hình FE $50\times50$ (đoạn tiếp theo).

\begin{table}[t]
\centering
\caption{Quy trình sinh (best-of-30 + tinh chỉnh nhận biết nhị phân hóa có
bảo vệ) so với tối ưu tô-pô chạy từ đầu (tốt nhất trong bốn điểm khởi tạo,
hội tụ) trên ba tập 100 mục tiêu không giao nhau. Cả hai phương pháp chọn
ứng viên theo sai số kết hợp kiểm chứng (không theo điểm tổng hợp ở
Bảng~\ref{tab:multi}). FE: số lần giải mỗi mục
tiêu. Chế tạo được: một pha rắn liên thông, nét tối thiểu hai pixel. Hai dòng
phân tầng: mức giảm tương đối sai số của quy trình sinh so với tối ưu tô-pô
cho mục tiêu $r<10$ (bootstrap ghép cặp, CI 95\%; dương nghiêng về quy trình
sinh). Quy trình cuối không ép biên tuần hoàn; kiểm định đăng ký trước trên
tập C chạy với bước này (xem văn bản; cả hai phiên bản trong Thông tin bổ
sung).}
\label{tab:simp}
\small
\setlength{\tabcolsep}{3.5pt}
\begin{tabular}{llccccc}
\toprule
Tập & Phương pháp & $\RR(\nuot)$ & $\RR(\nuto)$ & MAE kết hợp & FE & Chế tạo\\
\midrule
\multirow{2}{*}{A} & sinh & 0,998 & 0,846 & 0,0139 & 107 & 0,32\\
 & tối ưu tô-pô & 0,997 & 0,997 & 0,0128 & 897 & 0,77\\
\multirow{2}{*}{B} & sinh & 0,995 & 0,745 & 0,0229 & 107 & 0,35\\
 & tối ưu tô-pô & 0,996 & 0,997 & 0,0139 & 889 & 0,74\\
\multirow{2}{*}{C} & sinh & 0,998 & 0,998 & 0,0082 & 107 & 0,32\\
 & tối ưu tô-pô & 0,998 & 0,996 & 0,0126 & 893 & 0,83\\
\midrule
\multicolumn{7}{l}{$r<10$: sai số kết hợp \;A $+41\%$ [26; 54] \;B $+40\%$ [23; 53] \;C $+35\%$ [19; 48]}\\
\multicolumn{7}{l}{$r<10$: sai số $\nuot$ \;A $+39\%$ [19; 54] \;B $+39\%$ [19; 54] \;C $+37\%$ [14; 54]}\\
\bottomrule
\end{tabular}
\end{table}

\paragraph{Độ nhạy theo lưới.}
Mọi độ chính xác ở trên là so với mô hình FE $50\times50$ dùng cho sinh dữ
liệu, tinh chỉnh và kiểm chứng. Giải lại cùng hình học nhị phân với mỗi pixel
chia thành $4\times4$ phần tử làm $\nuot$ thay đổi trung vị 0,015 (3,8\%), tối
đa 0,042 trên 30 thiết kế test, theo hướng âm hơn ở 27 thiết kế. Trên các
thiết kế cuối của tập thứ nhất, thứ hạng hai phương pháp không đổi trên lưới
$200\times200$ ($\RR(\nuot)$ 0,992 so với 0,985).

\paragraph{Phương án lai không thừa hưởng điểm mạnh của cả hai.}
Vì tối ưu tô-pô bền vững hơn và dễ chế tạo hơn còn quy trình sinh rẻ hơn,
chúng tôi đăng ký trước một phương án lai trong đó thiết kế sinh ra làm điểm
khởi tạo cho tối ưu tô-pô (continuation $\beta=8\to64$, tối đa thêm 61 lần
giải FE). Phương án này đạt tiêu chí chi phí và không suy biến (166 lần giải
FE mỗi mục tiêu; 1\% suy biến) nhưng không đạt tiêu chí không-kém-hơn về sai
số kết hợp so với tối ưu tô-pô hội tụ (cận dưới CI $-96\%$, do mục tiêu duy
nhất bị đứt) và tiêu chí khả năng chế tạo (0,33 so với ngưỡng 0,70). Bắt đầu
với phép chiếu sắc giữ nguyên tô-pô, kể cả các nét mảnh, của thiết kế sinh ra.

% --- Phương pháp MỚI -------------------------------------------------------
\subsection{Baseline tối ưu tô-pô}
Với mỗi mục tiêu, chúng tôi cực tiểu
$\tfrac12\lVert\boldsymbol\nu(\hat{\mathbf x})-\boldsymbol\nu^{*}\rVert_2^2$
theo mật độ phần tử $\mathbf x$, với $\hat{\mathbf x}=H_\beta(\tilde{\mathbf x})$
và $\tilde{\mathbf x}$ là trường đã lọc mật độ (bán kính 1,5 phần tử), chịu
ràng buộc $\overline{\hat{\mathbf x}}\le0{,}55$ và
$Q_{11},Q_{22}\ge0{,}1\cdot0{,}55\,E_0$, bằng MMA \citep{svanberg1987mma} với
hệ số phạt $p=3$. Độ nhạy dùng cùng biểu thức giải tích như tinh chỉnh và đã
được kiểm tra bằng sai phân hữu hạn. Continuation chạy
$\beta\in\{1,2,4,8,16,32,64\}$, tối đa 43 lần đánh giá mỗi mức và khởi động
lại bộ tối ưu ở mỗi mức. Tô-pô khởi tạo là bốn họ seed phổ biến nhất trong dữ
liệu huấn luyện, co về cận thể tích. $\hat{\mathbf x}$ cuối được lấy ngưỡng 0,5
và kiểm chứng bằng một lần giải FE độc lập. Tham số là trung vị của dữ liệu
huấn luyện. Phương án lai dùng thiết kế nhị phân sinh ra làm điểm khởi tạo và
$\beta\in\{8,16,32,64\}$, tối đa 15 lần đánh giá mỗi mức.

% --- BỔ SUNG "Thống kê và đăng ký trước" ------------------------------------
Hai thí nghiệm nữa được ghi tiêu chí bằng văn bản trước khi chạy: phương án
lai (bốn tiêu chí: sai số kết hợp không kém hơn quá 10\%, ít nhất 70\% chế tạo
được, tối đa 299 lần giải FE mỗi mục tiêu, tối đa 1\% thiết kế suy biến) và so
sánh phân tầng trên tập mục tiêu thứ ba (cả hai CI không chứa 0 với $r<10$).
Hai tập mục tiêu đầu đã được phân tích trước khi nghĩ ra cách phân tầng, nên
kết quả phân tầng của chúng được báo cáo là hậu kiểm.

% --- THAY đoạn Giới hạn ----------------------------------------------------
\paragraph{Giới hạn.}
(i) Độ chính xác được đo so với mô hình FE $50\times50$ đàn hồi tuyến tính,
biến dạng nhỏ, với hệ số Poisson vật liệu nền 0,3; sai số rời rạc hóa của mô
hình này (trung vị 0,015 với $\nuot$) lớn hơn sai số trung bình được báo cáo.
Biến dạng lớn và ô ba chiều chưa được xét, và chưa thiết kế nào được thử
nghiệm thực tế.
(ii) Quy trình sinh thất bại với mục tiêu dị hướng mạnh ($r\ge10$), vốn thưa
trong dữ liệu huấn luyện; tối ưu tô-pô thì không.
(iii) Tối ưu tô-pô chạy từ đầu tạo ô chế tạo được gấp hơn hai lần
(0,74--0,83 so với 0,32--0,35). Tinh chỉnh không làm thay đổi điều này một
cách có hệ thống: qua năm lần chạy, tỉ lệ chế tạo được dịch chuyển từ $-6$
đến $+9$ điểm phần trăm (61 thiết kế trở nên chế tạo được, 44 thiết kế mất
tính chất này; kiểm định dấu $p=0{,}12$).
(iv) Với mục tiêu không auxetic ngoài dải huấn luyện, tối ưu tô-pô chính xác
hơn rõ rệt (MAE $\nuot$ 0,005--0,040 so với 0,06--0,09); lợi thế của mô hình
sinh ở vùng này chỉ đúng khi so với tra cứu thư viện huấn luyện.
(v) Lợi thế chi phí là lợi thế khấu hao: sinh dữ liệu huấn luyện cần khoảng
$2{,}7$--$3{,}0\times10^{6}$ lần giải FE, chỉ hòa vốn so với tối ưu tô-pô hội
tụ sau khoảng 3\,500 mục tiêu.
(vi) Điểm thẩm mỹ là một heuristic chưa được kiểm chứng với đánh giá của con
người.

% --- THAY các đoạn Discussion "Độ chính xác", "Đa mục tiêu", "Ngoài phân
% phối" (Giới hạn: xem khối trên). [C5] = chỗ đổi nếu bỏ ép biên tuần hoàn.
\paragraph{Tối ưu đúng thứ được kiểm chứng.}
Phát hiện trung tâm là mọi tìm kiếm dựa trên gradient trên trường mật độ, dù
trong không gian ẩn của mô hình sinh hay trên mật độ phần tử trong tối ưu
tô-pô, đều sẽ khai thác sự khác biệt giữa trường được lấy đạo hàm và thiết kế
nhị phân được phân tích sau cùng. Bộ tối ưu đạt hàm mục tiêu $10^{-10}$ trên
những trường mà bộ kiểm chứng nhìn thấy rất khác (Hình~\ref{fig:gap}), và điều
tương tự xảy ra trong tối ưu tô-pô với Heaviside continuation tới $\beta=64$.
Khoảng lệch này không hiện ra trên đường cong hàm mất mát. Có hai cách khắc
phục và cả hai đều rẻ: đưa các toán tử kiểm chứng vào hàm mục tiêu, và chọn
giữa các ứng viên theo sai số kiểm chứng thay vì sai số tối ưu. Cả hai thành
phần đều không mới trong tối ưu tô-pô dựa trên mật độ
\citep{guest2004achieving,wang2011projection}; điều công trình này chỉ ra là bỏ
qua chúng sẽ làm thay đổi kết luận khi so sánh các phương pháp, và con đường
mô hình sinh lẫn con đường kinh điển đều cần chúng như nhau.

\paragraph{Mô hình sinh đóng góp thêm điều gì.}
Khi đo so với tối ưu tô-pô hội tụ thay vì so với tra bảng, lợi thế của quy
trình sinh hẹp hơn so với điều thường được ngầm hiểu. Nó đạt cùng độ chính
xác kiểm chứng với số lần giải FE ít hơn khoảng tám lần mỗi mục tiêu, và sai
số kết hợp thấp hơn với mục tiêu dị hướng vừa phải, nhưng phần tiết kiệm chỉ
được khấu hao sau hàng nghìn mục tiêu vì chính dữ liệu huấn luyện được tạo
bằng tối ưu tô-pô. Vì vậy quy trình này hấp dẫn khi phải phục vụ nhiều mục
tiêu, chẳng hạn trong thiết kế tương tác hay trong vòng lặp trong của tối ưu
đa tỉ lệ, chứ không phải khi chỉ cần một ô. Các thất bại của nó dự đoán được:
mục tiêu dị hướng mạnh vốn thưa trong dữ liệu huấn luyện. Báo cáo các thất
bại này theo tầng, thay vì gộp vào một điểm trung bình, đã làm chúng lộ ra
trong nghiên cứu này và nên trở thành thông lệ.

\paragraph{Khả năng chế tạo và mục tiêu ngoài phân phối.}
Tối ưu tô-pô chạy từ đầu tạo ra ô liên thông với kích thước nét đủ lớn gấp
hơn hai lần so với quy trình sinh, và tinh chỉnh không gian ẩn không
làm thay đổi tỉ lệ này một cách có hệ thống. Khởi tạo tối ưu tô-pô bằng thiết kế sinh ra không khôi phục
được tính chất đó, vì phép chiếu sắc ban đầu giữ nguyên tô-pô được sinh. Với
mục tiêu không auxetic ngoài dải huấn luyện, cVAE đạt đúng dấu ở nơi tra cứu
thư viện không thể, nhưng tối ưu tô-pô đạt biên độ mục tiêu chính xác hơn
nhiều. Theo trục biên độ, mọi biến thể sinh đều dừng ở biên của dữ liệu huấn
luyện. Vì vậy, các khẳng định OOD cho mô hình sinh nên nêu rõ hướng trong
không gian tính chất, kiểm tra mục tiêu có khả thi về vật lý không, và so với
tối ưu tô-pô thay vì tra cứu. [C5: nếu bỏ ép biên tuần hoàn, thêm một câu rằng
điều kiện biên tuần hoàn dựa trên phần tử khiến mọi ô pixel đều lát được bằng
tịnh tiến, nên không cần khớp cạnh.]
